\chapter{xTB pre-screening run report (menu 42)}
\label{ch:menu42}
\menulabel{42}

\section{Purpose}

xTB (GFN2-xTB) is the Tier-1 pre-screening level of the protocol chain:
it decides what is worth a target-level calculation and seeds the geometry.
This menu reads a run capture back -- the task kinds, the energy, the
convergence markers, the frequency set and the thermochemistry summary.

\section{What you need}

The captured stdout of an xtb run (redirect both streams:
\code{xtb x.xyz -{}-gfn2 -{}-ohess > x.out 2>\&1}); the closing
\code{normal termination of xtb} line goes to stderr, so a stdout-only
capture lacks it.

\section{Walkthrough}

The example reads the water opt+hess capture -- converged after five
cycles, nine frequency values (six translation/rotation zeros plus the
three water modes), zero imaginary -- and then the planar-ammonia capture,
whose imaginary mode ($-1337.66$~cm$^{-1}$) is flagged with the engine's
own counter.

\lstinputlisting[style=fbkcode, firstline=57, lastline=79,
  title={The menu-42 example (the opt+hess capture: task kinds, the
  frequency set and the reading notes)}]
  {examples/expected/m42_xtb.txt}

\section{What you get}

The report \code{<capture>.xtb.fbk.md}: the program banner data, the task
kinds recognised (single point / geometry optimisation / frequencies),
the final energy, gradient norm and HOMO--LUMO gap, the SCF markers, the
optimisation outcome, the frequency set with the imaginary count
cross-checked against the engine's thermochemistry counters, and the
printed thermochemistry lines.

\section{Conventions, cross-checks and boundaries (measured)}

\begin{readbox}
\begin{itemize}
  \item \textbf{The frequency block is printed twice} (the Hessian step and
    the thermochemistry section); both copies are kept as measured, and
    the report counts imaginary modes from the first group.
  \item \textbf{Both optimisation markers are read}: ``CONVERGED AFTER N
    ITERATIONS'' and -- when \code{-{}-cycles} ran out -- ``FAILED TO
    CONVERGE \ldots IN N ITERATIONS''; the latter is reported with the
    warning that the geometry is a partial step.
  \item \textbf{The imaginary count is cross-checked} against the engine's
    \code{\# imaginary freq.} counter and the ``found N significant
    imaginary frequency'' line (the engine's own 5~cm$^{-1}$ cut-off); an
    imaginary mode marks a saddle or a non-stationary geometry at this
    level and must be re-verified at the target level.
  \item \textbf{Pre-screening discipline}: GFN2 structures and energies
    rank and seed, they are not carried upward as results (measured
    example: a GFN2 transition state re-optimised at r2SCAN-3c carried
    seven imaginary modes).  GFN2 covers the fifteen lanthanides; the
    actinides have no semi-empirical coverage in this stack.
  \item \textbf{A measured non-feature}: linear water gives five zero
    modes and a positive degenerate bend in this engine's harmonic
    picture -- the report relays what the file carries, without physical
    extrapolation.
\end{itemize}
\end{readbox}
